% This document was generated by an LLM.

\documentclass{workbook}

\usepackage{amsthm}

\theoremstyle{definition}
\newtheorem{exercise}{Exercise}[section]

\newcommand{\N}{\mathbb{N}}
\newcommand{\Z}{\mathbb{Z}}
\newcommand{\Q}{\mathbb{Q}}
\newcommand{\R}{\mathbb{R}}
\newcommand{\ans}[1]{\paragraph{Exercise~\ref{#1}.}}

\title{Negating Definitions\\[0.3em]\large A Workbook on Quantifiers in Real Analysis}
\author{}
\date{}

\begin{document}
\maketitle

%======================================================================
\section{How to use this workbook}

Each exercise gives a statement. Unless the exercise says otherwise, do the following.
\begin{enumerate}
  \item Write the statement with explicit quantifiers, if it is not already written that way.
  \item Negate it mechanically: push $\neg$ inward until it stands only in front of atomic statements, then eliminate it.
  \item Translate the negation into one plain sentence.
  \item Carry out any test the exercise asks for. When a negation begins with $\exists$, construct the witness explicitly; when a witness must be chosen after a $\forall$, write it as a function of the variables chosen before it.
\end{enumerate}
Answers are in Section~\ref{sec:answers}.

\textbf{Conventions.} $\N = \{1, 2, 3, \dots\}$. Unless stated otherwise, $A \subseteq \R$, $f \colon \R \to \R$, and $(a_n)$ is a real sequence indexed by $\N$.

%======================================================================
\section{The rules}

\begin{align*}
  \neg\, \forall x\ P(x) &\equiv \exists x\ \neg P(x) &
  \neg\, \exists x\ P(x) &\equiv \forall x\ \neg P(x) \\
  \neg (P \land Q) &\equiv \neg P \lor \neg Q &
  \neg (P \lor Q) &\equiv \neg P \land \neg Q \\
  \neg (P \implies Q) &\equiv P \land \neg Q &
  \neg (P \iff Q) &\equiv (P \land \neg Q) \lor (Q \land \neg P) \\
  \neg (a < b) &\equiv a \ge b &
  \neg (a \le b) &\equiv a > b
\end{align*}

\textbf{Bounded quantifiers.} The notation $\forall x \in A : P$ abbreviates $\forall x\ (x \in A \implies P)$, and $\exists x \in A : P$ abbreviates $\exists x\ (x \in A \land P)$. Applying the rules above gives
\[
  \neg (\forall x \in A : P) \equiv \exists x \in A : \neg P,
  \qquad
  \neg (\exists x \in A : P) \equiv \forall x \in A : \neg P.
\]
The restriction survives negation. Negating $\forall \epsilon > 0$ gives $\exists \epsilon > 0$, never $\epsilon \le 0$. The same holds for $\forall n \ge N$, $\forall x \in D$, and so on.

\textbf{Common errors.}
\begin{itemize}
  \item Negating the restriction of a bounded quantifier.
  \item Writing the contrapositive $\neg Q \implies \neg P$, which is equivalent to $P \implies Q$, not its negation. Writing the converse $Q \implies P$, which is unrelated to it.
  \item Replacing a word by its ``opposite.'' Not open is not closed; not increasing is not decreasing; not even is not odd.
  \item Forgetting that negation swaps strict and non-strict inequalities.
\end{itemize}

%======================================================================
\section{Warm-up: one quantifier}

\begin{exercise}\label{ex:pos}
Every element of $A$ is positive.
\end{exercise}

\begin{exercise}\label{ex:int}
$A$ contains an integer.
\end{exercise}

\begin{exercise}\label{ex:zero}
$f$ has a zero. Test: show that $f(x) = e^x$ satisfies the negation.
\end{exercise}

\begin{exercise}\label{ex:nonneg}
$f$ is nonnegative. Test: show that $\sin$ satisfies the negation.
\end{exercise}

\begin{exercise}\label{ex:even}
$f$ is even, i.e.\ $\forall x \in \R : f(-x) = f(x)$. Test: show that $f(x) = x^2 + x$ is not even. Is it odd?
\end{exercise}

\begin{exercise}\label{ex:vacuous}
Take $A = \emptyset$. Negate ``every element of $A$ is positive.'' Decide which of the statement and its negation is true.
\end{exercise}

\begin{exercise}\label{ex:impl}
Negate $\forall x \in \R : (x > 1 \implies f(x) > 0)$. Then classify each of the following as the negation, an equivalent statement, or neither:
\begin{enumerate}
  \item $\forall x \in \R : (x > 1 \implies f(x) \le 0)$;
  \item $\forall x \in \R : (f(x) \le 0 \implies x \le 1)$;
  \item $\forall x \in \R : (f(x) > 0 \implies x > 1)$.
\end{enumerate}
\end{exercise}

%======================================================================
\section{Two quantifiers}

\begin{exercise}\label{ex:bddabove}
$A$ is bounded above: $\exists M \in \R\ \forall a \in A : a \le M$. Test: show that $\N$ is not bounded above.
\end{exercise}

\begin{exercise}\label{ex:bdd}
$A$ is bounded: $\exists M \in \R\ \forall a \in A : |a| \le M$.
\end{exercise}

\begin{exercise}\label{ex:max}
$A$ has a maximum: $\exists m \in A\ \forall a \in A : a \le m$. Test: show that $(0,1)$ has no maximum. Does $\emptyset$ have one?
\end{exercise}

\begin{exercise}\label{ex:inj}
$f$ is injective: $\forall x, y \in \R : (f(x) = f(y) \implies x = y)$. Test: $f(x) = x^2$.
\end{exercise}

\begin{exercise}\label{ex:surj}
$f$ is surjective: $\forall y \in \R\ \exists x \in \R : f(x) = y$. Test: $f(x) = e^x$.
\end{exercise}

\begin{exercise}\label{ex:periodic}
$f$ is periodic: $\exists p > 0\ \forall x \in \R : f(x + p) = f(x)$. Test: show that $f(x) = x^2$ is not periodic.
\end{exercise}

\begin{exercise}\label{ex:swap1}
Over $\R$, compare
\[
  \text{(i) } \forall x\ \exists y : y > x
  \qquad\text{and}\qquad
  \text{(ii) } \exists y\ \forall x : y > x.
\]
Translate both. Decide which is true, and prove the false one false by proving its negation.
\end{exercise}

\begin{exercise}\label{ex:incr}
$f$ is increasing: $\forall x, y \in \R : (x < y \implies f(x) \le f(y))$. Test: $f(x) = x^2$.
\end{exercise}

\begin{exercise}\label{ex:monotone}
$f$ is monotone, meaning increasing or decreasing. Write the full quantified statement, negate it, and show that $\sin$ satisfies the negation. How many witnesses does the negation require?
\end{exercise}

\begin{exercise}\label{ex:const}
$f$ is constant. Write it in two ways:
\[
  \text{(i) } \exists c \in \R\ \forall x \in \R : f(x) = c
  \qquad\text{and}\qquad
  \text{(ii) } \forall x, y \in \R : f(x) = f(y).
\]
Negate both and prove the two negations equivalent.
\end{exercise}

%======================================================================
\section{Sequences: three quantifiers}

\begin{exercise}\label{ex:eventually}
``$a_n > 0$ for all sufficiently large $n$'' is formalized as $\exists N \in \N\ \forall n \ge N : a_n > 0$. Negate it. Then prove that for any property $P$ of natural numbers,
\[
  \forall N \in \N\ \exists n \ge N : P(n)
  \iff
  \{ n \in \N : P(n) \} \text{ is infinite.}
\]
\end{exercise}

\begin{exercise}\label{ex:conv}
$a_n \to L$:
\[
  \forall \epsilon > 0\ \exists N \in \N\ \forall n \ge N : |a_n - L| < \epsilon.
\]
Test: show that $a_n = (-1)^n$ does not converge to $1$.
\end{exercise}

\begin{exercise}\label{ex:div}
$(a_n)$ converges: $\exists L \in \R : a_n \to L$. Write the negation (``$(a_n)$ diverges'') in full. Test: show that $(-1)^n$ diverges.
\end{exercise}

\begin{exercise}\label{ex:cauchy}
$(a_n)$ is Cauchy:
\[
  \forall \epsilon > 0\ \exists N \in \N\ \forall m, n \ge N : |a_m - a_n| < \epsilon.
\]
Test: show that $H_n = \sum_{k=1}^{n} 1/k$ is not Cauchy.
\end{exercise}

\begin{exercise}\label{ex:infty}
$a_n \to \infty$: $\forall M \in \R\ \exists N \in \N\ \forall n \ge N : a_n > M$.
\end{exercise}

\begin{exercise}\label{ex:unbdd}
Compare ``$(a_n)$ is unbounded above'' ($\forall M\ \exists n : a_n > M$) with $a_n \to \infty$.
\begin{enumerate}
  \item Prove that $a_n \to \infty$ implies $(a_n)$ is unbounded above.
  \item Let $a_n = n$ for even $n$ and $a_n = 0$ for odd $n$. Show it is unbounded above and that $a_n \not\to \infty$.
\end{enumerate}
\end{exercise}

\begin{exercise}\label{ex:swapconv}
Move $\exists N$ to the front of the definition of convergence:
\[
  \exists N \in \N\ \forall \epsilon > 0\ \forall n \ge N : |a_n - L| < \epsilon.
\]
What does this say about $(a_n)$? Prove your answer.
\end{exercise}

\begin{exercise}\label{ex:cluster}
Change $\exists N\ \forall n \ge N$ to $\forall N\ \exists n \ge N$:
\[
  \forall \epsilon > 0\ \forall N \in \N\ \exists n \ge N : |a_n - L| < \epsilon.
\]
Translate it using Exercise~\ref{ex:eventually}. Show that $(-1)^n$ satisfies it with $L = 1$, so it is strictly weaker than convergence.
\end{exercise}

%======================================================================
\section{Suprema}

Throughout this section $A \neq \emptyset$.

\begin{exercise}\label{ex:ub}
$u$ is an upper bound of $A$.
\end{exercise}

\begin{exercise}\label{ex:sup}
$s = \sup A$ is characterized by
\[
  (\forall a \in A : a \le s) \land (\forall \epsilon > 0\ \exists a \in A : a > s - \epsilon).
\]
Negate it. Test: prove that $1 = \sup (0,1)$. Show that $2 \ne \sup [0,1]$ and $\tfrac12 \ne \sup (0,1)$, and say which clause fails in each case.
\end{exercise}

\begin{exercise}\label{ex:lub}
The textbook definition: $s = \sup A$ if $s$ is an upper bound of $A$ and $s \le u$ for every upper bound $u$ of $A$. Expand ``$u$ is an upper bound'' into quantifiers, so that the second clause becomes
\[
  \forall u \in \R : \bigl( (\forall a \in A : a \le u) \implies s \le u \bigr).
\]
Negate the second clause. Which quantifier is negated and which is not?
\end{exercise}

%======================================================================
\section{Limits and continuity}

In this section $f \colon D \to \R$ with $D \subseteq \R$.

\begin{exercise}\label{ex:lim}
$\lim_{x \to c} f(x) = L$, for $c$ a limit point of $D$:
\[
  \forall \epsilon > 0\ \exists \delta > 0\ \forall x \in D : \bigl(0 < |x - c| < \delta \implies |f(x) - L| < \epsilon\bigr).
\]
Test: let $f(x) = \operatorname{sgn}(x)$ on $D = \R \setminus \{0\}$. Show that $\lim_{x \to 0} f(x) \ne L$ for every $L \in \R$.
\end{exercise}

\begin{exercise}\label{ex:cont}
$f$ is continuous at $c \in D$:
\[
  \forall \epsilon > 0\ \exists \delta > 0\ \forall x \in D : \bigl(|x - c| < \delta \implies |f(x) - f(c)| < \epsilon\bigr).
\]
Test: Dirichlet's function ($1$ on $\Q$, $0$ elsewhere) is discontinuous at every $c \in \R$.
\end{exercise}

\begin{exercise}\label{ex:isolated}
Let $D = \{0\} \cup [1, 2]$ and let $f \colon D \to \R$ be arbitrary. Prove that $f$ is continuous at $0$. Explain why every attempt to satisfy the negation fails.
\end{exercise}

\begin{exercise}\label{ex:xq}
Let $f(x) = x$ for $x \in \Q$ and $f(x) = 0$ otherwise. Prove that $f$ is continuous at $0$ and discontinuous at every $c \ne 0$.
\end{exercise}

\begin{exercise}\label{ex:uc}
$f$ is uniformly continuous on $D$:
\[
  \forall \epsilon > 0\ \exists \delta > 0\ \forall x, y \in D : \bigl(|x - y| < \delta \implies |f(x) - f(y)| < \epsilon\bigr).
\]
Test: show that $f(x) = x^2$ is not uniformly continuous on $\R$.
\end{exercise}

\begin{exercise}\label{ex:ucrecip}
Show that $f(x) = 1/x$ is not uniformly continuous on $(0,1)$.
\end{exercise}

\begin{exercise}\label{ex:ucimpl}
Write ``$f$ is continuous on $D$'' (continuous at every $c \in D$) with all quantifiers. Compare it with uniform continuity, identify the quantifier move, and prove that uniform continuity implies continuity. Which earlier exercises show the converse fails?
\end{exercise}

\begin{exercise}\label{ex:lip}
$f$ is Lipschitz on $D$: $\exists K \ge 0\ \forall x, y \in D : |f(x) - f(y)| \le K |x - y|$. Test: show that $\sqrt{x}$ is not Lipschitz on $[0,1]$, but is uniformly continuous there.
\end{exercise}

%======================================================================
\section{Sequences of functions}

Let $f_n, f \colon D \to \R$.

\begin{exercise}\label{ex:pw}
Pointwise and uniform convergence:
\begin{align*}
  f_n \to f \text{ pointwise} &\iff \forall x \in D\ \forall \epsilon > 0\ \exists N\ \forall n \ge N : |f_n(x) - f(x)| < \epsilon, \\
  f_n \to f \text{ uniformly} &\iff \forall \epsilon > 0\ \exists N\ \forall n \ge N\ \forall x \in D : |f_n(x) - f(x)| < \epsilon.
\end{align*}
Negate both. Identify the quantifier move that separates them and say what $N$ may depend on in each.
\end{exercise}

\begin{exercise}\label{ex:xn}
Let $f_n(x) = x^n$ on $D = [0,1)$. Show $f_n \to 0$ pointwise but not uniformly.
\end{exercise}

\begin{exercise}\label{ex:unifdomain}
\begin{enumerate}
  \item Show that $x^n \to 0$ uniformly on $[0, a]$ for every $0 < a < 1$.
  \item Let $f_n(x) = x/n$. Show $f_n \to 0$ pointwise on $\R$ but not uniformly, and uniformly on $[-R, R]$ for every $R > 0$.
\end{enumerate}
\end{exercise}

%======================================================================
\section{Sets}

\begin{exercise}\label{ex:dense}
$A$ is dense in $\R$: $\forall a, b \in \R : \bigl(a < b \implies \exists q \in A : a < q < b\bigr)$. Test: show that $\Z$ is not dense and that $\Q$ is.
\end{exercise}

\begin{exercise}\label{ex:open}
$U \subseteq \R$ is open:
\[
  \forall x \in U\ \exists r > 0\ \forall y \in \R : \bigl(|y - x| < r \implies y \in U\bigr).
\]
A set is closed if its complement is open. Show that $[0,1)$ is neither open nor closed, and that $\emptyset$ is both.
\end{exercise}

\begin{exercise}\label{ex:limpt}
$c$ is a limit point of $A$: $\forall \delta > 0\ \exists a \in A : 0 < |a - c| < \delta$. A point $c \in A$ that is not a limit point of $A$ is isolated. Show that every point of $\Z$ is isolated in $\Z$, and that $0$ is a limit point of $\{1/n : n \in \N\}$ not belonging to it.
\end{exercise}

\begin{exercise}\label{ex:seqcpt}
$K \subseteq \R$ is sequentially compact: every sequence in $K$ has a subsequence converging to a point of $K$. Write this with quantifiers over sequences, subsequences and points, and negate it. Show that $(0,1]$ and $\R$ are not sequentially compact, and explain why they fail for different reasons.
\end{exercise}

%======================================================================
\section{Synthesis}

\begin{exercise}\label{ex:equi}
A family $\mathcal{F}$ of functions $D \to \R$ is equicontinuous on $D$ if
\[
  \forall \epsilon > 0\ \exists \delta > 0\ \forall f \in \mathcal{F}\ \forall x, y \in D : \bigl(|x - y| < \delta \implies |f(x) - f(y)| < \epsilon\bigr).
\]
Negate it. Show that $\{x^n : n \in \N\}$ is not equicontinuous on $[0,1]$, even though each member is uniformly continuous there.
\end{exercise}

\begin{exercise}\label{ex:everycont}
Write ``every continuous function on $D$ is bounded'' with quantifiers, including a quantifier over functions. Negate it. Show that the negation holds for $D = (0,1)$. What does the extreme value theorem say about $D = [0,1]$?
\end{exercise}

\begin{exercise}\label{ex:swapgeneral}
For an arbitrary predicate $P(x, y)$, prove that
\[
  \exists y\ \forall x\ P(x,y) \implies \forall x\ \exists y\ P(x,y),
\]
and give a counterexample to the converse. Show that two adjacent quantifiers of the same kind always commute. Which exercises in this workbook are instances of these facts?
\end{exercise}

%======================================================================
\newpage
\section{Answers}\label{sec:answers}

\subsection*{Warm-up}

\ans{ex:pos}
$\forall a \in A : a > 0$. Negation: $\exists a \in A : a \le 0$. Some element of $A$ is not positive.

\ans{ex:int}
$\exists a \in A : a \in \Z$. Negation: $\forall a \in A : a \notin \Z$. No element of $A$ is an integer.

\ans{ex:zero}
$\exists x \in \R : f(x) = 0$. Negation: $\forall x \in \R : f(x) \ne 0$. For $e^x$: $e^x > 0$ for all $x$.

\ans{ex:nonneg}
$\forall x : f(x) \ge 0$. Negation: $\exists x : f(x) < 0$. Witness for $\sin$: $x = -\pi/2$, $\sin(-\pi/2) = -1$.

\ans{ex:even}
Negation: $\exists x : f(-x) \ne f(x)$. Witness: $x = 1$, since $f(-1) = 0 \ne 2 = f(1)$. The function is not odd either: oddness is $\forall x : f(-x) = -f(x)$, and $f(-1) = 0 \ne -2 = -f(1)$. ``Not even'' and ``odd'' are different statements.

\ans{ex:vacuous}
Negation: $\exists a \in \emptyset : a \le 0$. This is false, because nothing is in $\emptyset$. Hence the original statement is true. Every universal statement over the empty set is true.

\ans{ex:impl}
Negation: $\exists x : (x > 1 \land f(x) \le 0)$. Statement 1 is neither: for $f = \sin$, both the original and statement 1 are false, so statement 1 is not the negation. Statement 2 is the contrapositive, hence equivalent. Statement 3 is the converse, so neither: $f(x) = 1$ satisfies the original but not the converse.

\subsection*{Two quantifiers}

\ans{ex:bddabove}
Negation: $\forall M \in \R\ \exists a \in A : a > M$. For $A = \N$: given $M$, the Archimedean property gives $n \in \N$ with $n > M$.

\ans{ex:bdd}
Negation: $\forall M \in \R\ \exists a \in A : |a| > M$. Some element exceeds every proposed bound in absolute value.

\ans{ex:max}
Negation: $\forall m \in A\ \exists a \in A : a > m$. For $(0,1)$: given $m \in (0,1)$, take $a = (m+1)/2$. Then $a \in (0,1)$ and $a > m$. For $\emptyset$, the negation $\forall m \in \emptyset\ \dots$ is vacuously true, so $\emptyset$ has no maximum.

\ans{ex:inj}
Negation: $\exists x, y : (f(x) = f(y) \land x \ne y)$. For $x^2$: $x = 1$, $y = -1$.

\ans{ex:surj}
Negation: $\exists y\ \forall x : f(x) \ne y$. Some value is never attained. For $e^x$: $y = 0$.

\ans{ex:periodic}
Negation: $\forall p > 0\ \exists x : f(x + p) \ne f(x)$. For $x^2$: given $p > 0$, take $x = 0$; then $f(p) = p^2 \ne 0 = f(0)$. The witness $x$ happens not to depend on $p$ here; in general it may.

\ans{ex:swap1}
(i) says there is no largest real number; it is true, with $y = x + 1$. (ii) says some real number exceeds every real number; it is false. Negation of (ii): $\forall y\ \exists x : x \ge y$. Given $y$, take $x = y$.

\ans{ex:incr}
Negation: $\exists x, y : (x < y \land f(x) > f(y))$. For $x^2$: $x = -1$, $y = 0$.

\ans{ex:monotone}
Monotone:
\[
  \bigl(\forall x, y : x < y \implies f(x) \le f(y)\bigr) \lor \bigl(\forall x, y : x < y \implies f(x) \ge f(y)\bigr).
\]
Negation:
\[
  \bigl(\exists x_1, y_1 : x_1 < y_1 \land f(x_1) > f(y_1)\bigr) \land \bigl(\exists x_2, y_2 : x_2 < y_2 \land f(x_2) < f(y_2)\bigr).
\]
Two pairs are required, one for each clause. For $\sin$: the pair $(\pi/2, \pi)$ gives $1 > 0$, so $\sin$ is not increasing; the pair $(0, \pi/2)$ gives $0 < 1$, so it is not decreasing.

\ans{ex:const}
Negation of (i): $\forall c\ \exists x : f(x) \ne c$. Negation of (ii): $\exists x, y : f(x) \ne f(y)$. Equivalence: if $f(x) \ne f(y)$, then for any $c$, at least one of $f(x)$, $f(y)$ differs from $c$, which gives the required witness. Conversely, apply the negation of (i) with $c = f(0)$ to get $x$ with $f(x) \ne f(0)$; then the pair $(x, 0)$ works.

\subsection*{Sequences}

\ans{ex:eventually}
Negation: $\forall N\ \exists n \ge N : a_n \le 0$. That is, $a_n \le 0$ for infinitely many $n$.

Proof of the equivalence. Let $S = \{n : P(n)\}$. If $S$ is infinite and $N$ is given, then $S \not\subseteq \{1, \dots, N-1\}$, so some $n \in S$ has $n \ge N$. If $S$ is finite, take $N = 1 + \max S$ (or $N = 1$ if $S = \emptyset$); then no $n \ge N$ lies in $S$, so the left side fails.

In short: $\exists N\ \forall n \ge N$ means ``eventually,'' $\forall N\ \exists n \ge N$ means ``infinitely often,'' and they are each other's duals under negation.

\ans{ex:conv}
Negation: $\exists \epsilon > 0\ \forall N\ \exists n \ge N : |a_n - L| \ge \epsilon$. Some $\epsilon$-band around $L$ is left infinitely often. For $(-1)^n$ and $L = 1$: take $\epsilon = 1$. Given $N$, let $n \ge N$ be odd; then $|a_n - 1| = 2 \ge 1$.

\ans{ex:div}
Diverges:
\[
  \forall L \in \R\ \exists \epsilon > 0\ \forall N\ \exists n \ge N : |a_n - L| \ge \epsilon.
\]
For $(-1)^n$: fix $L$ and take $\epsilon = 1$. Since $|1 - L| + |{-1} - L| \ge |1 - (-1)| = 2$, at least one of $|1 - L|$, $|{-1} - L|$ is at least $1$. Both values $1$ and $-1$ occur for infinitely many $n$, so for every $N$ there is $n \ge N$ with $a_n$ equal to the value that is at distance at least $1$ from $L$. Note that $\epsilon$ is chosen after $L$; here the same $\epsilon$ happens to work for every $L$.

\ans{ex:cauchy}
Negation: $\exists \epsilon > 0\ \forall N\ \exists m, n \ge N : |a_m - a_n| \ge \epsilon$. For $H_n$: take $\epsilon = 1/2$. Given $N$, take $n = N$ and $m = 2N$. Then
\[
  H_{2N} - H_N = \sum_{k=N+1}^{2N} \frac{1}{k} \ge N \cdot \frac{1}{2N} = \frac12.
\]

\ans{ex:infty}
Negation: $\exists M\ \forall N\ \exists n \ge N : a_n \le M$. Some level $M$ is reached or undercut infinitely often.

\ans{ex:unbdd}
\begin{enumerate}
  \item Given $M$, the definition of $a_n \to \infty$ gives $N$ with $a_n > M$ for all $n \ge N$. In particular $a_N > M$.
  \item Unbounded: given $M$, take an even $n > M$; then $a_n = n > M$. Not $\to \infty$: use the negation from Exercise~\ref{ex:infty} with $M = 0$. Given $N$, take an odd $n \ge N$; then $a_n = 0 \le 0$.
\end{enumerate}

\ans{ex:swapconv}
It says that $a_n = L$ for all $n \ge N$: the sequence is eventually constant equal to $L$. Proof: for $n \ge N$, the number $|a_n - L| \ge 0$ is less than every $\epsilon > 0$, so it equals $0$. Conversely, an eventually constant sequence satisfies the statement. Moving $\exists N$ before $\forall \epsilon$ forbids $N$ from depending on $\epsilon$, and that destroys the notion of approximation.

\ans{ex:cluster}
By Exercise~\ref{ex:eventually}: for every $\epsilon > 0$, $|a_n - L| < \epsilon$ for infinitely many $n$. That is, $L$ is a cluster point (subsequential limit) of $(a_n)$. For $(-1)^n$ with $L = 1$: every even $n$ gives $|a_n - 1| = 0 < \epsilon$. Since $(-1)^n$ diverges (Exercise~\ref{ex:div}), the statement is strictly weaker than convergence.

\subsection*{Suprema}

\ans{ex:ub}
$\forall a \in A : a \le u$. Negation: $\exists a \in A : a > u$.

\ans{ex:sup}
Negation:
\[
  (\exists a \in A : a > s) \lor (\exists \epsilon > 0\ \forall a \in A : a \le s - \epsilon).
\]
Either $s$ is not an upper bound, or $s$ is an upper bound that can be lowered by some $\epsilon$.

$1 = \sup (0,1)$: every $a \in (0,1)$ satisfies $a < 1$. Given $\epsilon > 0$, let $a = \max(1 - \epsilon/2,\ 1/2)$. Then $1/2 \le a < 1$, so $a \in (0,1)$, and $a \ge 1 - \epsilon/2 > 1 - \epsilon$.

$2 \ne \sup [0,1]$: the second clause fails with $\epsilon = 1$, since every $a \in [0,1]$ satisfies $a \le 1 = 2 - \epsilon$.

$\tfrac12 \ne \sup (0,1)$: the first clause fails, with witness $a = 3/4$.

\ans{ex:lub}
Negation of the second clause:
\[
  \exists u \in \R : \bigl( (\forall a \in A : a \le u) \land s > u \bigr).
\]
Some upper bound lies strictly below $s$. The outer $\forall u$ is negated. The inner $\forall a$ is not, because it sits inside the antecedent of the implication, and $\neg(P \implies Q) \equiv P \land \neg Q$ leaves $P$ intact. Negation acts on the outermost structure only; scope determines which quantifiers flip.

\subsection*{Limits and continuity}

\ans{ex:lim}
Negation:
\[
  \exists \epsilon > 0\ \forall \delta > 0\ \exists x \in D : \bigl(0 < |x - c| < \delta \land |f(x) - L| \ge \epsilon\bigr).
\]
For $\operatorname{sgn}$ at $0$: fix $L$. Since $|1 - L| + |{-1} - L| \ge 2$, choose $\sigma \in \{1, -1\}$ with $|\sigma - L| \ge 1$. Take $\epsilon = 1$. Given $\delta > 0$, let $x = \sigma \delta / 2$. Then $0 < |x| < \delta$ and $|f(x) - L| = |\sigma - L| \ge 1$.

\ans{ex:cont}
Negation:
\[
  \exists \epsilon > 0\ \forall \delta > 0\ \exists x \in D : \bigl(|x - c| < \delta \land |f(x) - f(c)| \ge \epsilon\bigr).
\]
For Dirichlet's function: take $\epsilon = 1/2$. Given $\delta$, if $c \in \Q$ choose an irrational $x$ with $|x - c| < \delta$; if $c \notin \Q$ choose a rational one. Both exist by density of $\Q$ and of $\R \setminus \Q$. In either case $|f(x) - f(c)| = 1 \ge 1/2$.

\ans{ex:isolated}
Given $\epsilon > 0$, take $\delta = 1$. The only $x \in D$ with $|x - 0| < 1$ is $x = 0$, and $|f(0) - f(0)| = 0 < \epsilon$. The negation requires, for every $\delta$, some $x \in D$ with $|x| < \delta$ and $|f(x) - f(0)| \ge \epsilon > 0$. For $\delta \le 1$, the only candidate is $x = 0$, which gives $0 \ge \epsilon$, a contradiction. The neighborhood contains no other points of $D$ that could serve as witnesses.

\ans{ex:xq}
At $0$: $|f(x) - f(0)| = |f(x)| \le |x|$ for all $x$, so $\delta = \epsilon$ works.

At $c \ne 0$: take $\epsilon = |c|/2$ and let $\delta > 0$ be given. Put $\rho = \min(\delta, |c|/2)$.
\begin{itemize}
  \item If $c \in \Q$, then $f(c) = c$. Choose an irrational $x$ with $|x - c| < \rho$. Then $f(x) = 0$ and $|f(x) - f(c)| = |c| \ge \epsilon$.
  \item If $c \notin \Q$, then $f(c) = 0$. Choose a rational $x$ with $|x - c| < \rho$. Then $|x| > |c| - |c|/2 = |c|/2$, so $|f(x) - f(c)| = |x| > \epsilon$.
\end{itemize}
In both cases $|x - c| < \delta$. The restriction to $\rho \le |c|/2$ keeps $x$ away from $0$, where $f$ is small.

\ans{ex:uc}
Negation:
\[
  \exists \epsilon > 0\ \forall \delta > 0\ \exists x, y \in D : \bigl(|x - y| < \delta \land |f(x) - f(y)| \ge \epsilon\bigr).
\]
For $x^2$ on $\R$: take $\epsilon = 1$. Given $\delta$, let $x = 1/\delta$ and $y = x + \delta/2$. Then $|x - y| = \delta/2 < \delta$ and
\[
  y^2 - x^2 = (y - x)(y + x) = \frac{\delta}{2}\Bigl(\frac{2}{\delta} + \frac{\delta}{2}\Bigr) = 1 + \frac{\delta^2}{4} \ge 1.
\]
The witnesses move to infinity as $\delta \to 0$, where the slope of $x^2$ is unbounded.

\ans{ex:ucrecip}
Take $\epsilon = 1$. Given $\delta$, let $t = \min(\delta, 1/2)$, $x = t$, $y = t/2$. Both lie in $(0,1)$, $|x - y| = t/2 < \delta$, and
\[
  \frac{1}{y} - \frac{1}{x} = \frac{2}{t} - \frac{1}{t} = \frac{1}{t} \ge 2 \ge 1.
\]
Here the witnesses move toward $0$, the boundary point where $1/x$ blows up.

\ans{ex:ucimpl}
Continuous on $D$:
\[
  \forall c \in D\ \forall \epsilon > 0\ \exists \delta > 0\ \forall x \in D : \bigl(|x - c| < \delta \implies |f(x) - f(c)| < \epsilon\bigr).
\]
Reordering $\forall c\ \forall \epsilon$ to $\forall \epsilon\ \forall c$ changes nothing, since same-type quantifiers commute. Uniform continuity then moves $\exists \delta$ in front of $\forall c$ (renamed $x$, with the old $x$ renamed $y$). So $\delta$ may depend on $c$ in continuity but not in uniform continuity. Implication: given $\epsilon$ and $c$, the uniform $\delta$ for $\epsilon$ works at $c$ by taking $x = c$. Exercises~\ref{ex:uc} and~\ref{ex:ucrecip} give continuous functions that are not uniformly continuous.

\ans{ex:lip}
Negation: $\forall K \ge 0\ \exists x, y \in D : |f(x) - f(y)| > K |x - y|$. For $\sqrt{x}$ on $[0,1]$: given $K$, take $y = 0$ and $x = 1/(K+1)^2$. Then
\[
  \sqrt{x} - \sqrt{0} = \frac{1}{K+1} > \frac{K}{(K+1)^2} = K |x - 0|.
\]
Uniform continuity: for $x, y \ge 0$, $|\sqrt{x} - \sqrt{y}| \le \sqrt{|x - y|}$. To see this, assume $x \ge y$; then $\sqrt{x} \le \sqrt{y} + \sqrt{x - y}$ because squaring the right side gives $x + 2\sqrt{y(x-y)} \ge x$. Hence $\delta = \epsilon^2$ works. So Lipschitz is strictly stronger than uniform continuity.

\subsection*{Sequences of functions}

\ans{ex:pw}
Not pointwise:
\[
  \exists x \in D\ \exists \epsilon > 0\ \forall N\ \exists n \ge N : |f_n(x) - f(x)| \ge \epsilon.
\]
Not uniform:
\[
  \exists \epsilon > 0\ \forall N\ \exists n \ge N\ \exists x \in D : |f_n(x) - f(x)| \ge \epsilon.
\]
Reorder pointwise as $\forall \epsilon\ \forall x\ \exists N$; uniform moves $\forall x$ to after $\exists N$. In pointwise convergence $N$ may depend on $\epsilon$ and $x$; in uniform convergence on $\epsilon$ only. In the negation of uniform convergence, the bad point $x$ may change with $n$; in the negation of pointwise convergence, it is a single fixed point.

\ans{ex:xn}
Pointwise: for $x = 0$, $f_n(0) = 0$. For $0 < x < 1$, $x^n \to 0$; explicitly, given $\epsilon$, choose $N$ with $x^N < \epsilon$, which exists since $x^n \to 0$, and use $x^n \le x^N$ for $n \ge N$.

Not uniform: take $\epsilon = 1/2$. Given $N$, let $n = N$ and $x = 2^{-1/N}$. Then $x \in (0,1)$ and $x^N = 1/2 \ge \epsilon$.

\ans{ex:unifdomain}
\begin{enumerate}
  \item For $x \in [0, a]$: $|x^n| \le a^n$. Given $\epsilon$, choose $N$ with $a^N < \epsilon$. Then for $n \ge N$ and all $x \in [0,a]$, $x^n \le a^n \le a^N < \epsilon$. The bound $a^n$ does not involve $x$; that is what makes the convergence uniform.
  \item Pointwise: $|x/n| < \epsilon$ once $n > |x|/\epsilon$. Not uniform on $\R$: take $\epsilon = 1$; given $N$, let $n = N$ and $x = N$, so $|f_n(x)| = 1$. Uniform on $[-R, R]$: $|x/n| \le R/n < \epsilon$ for all $n > R/\epsilon$, independently of $x$.
\end{enumerate}

\subsection*{Sets}

\ans{ex:dense}
Negation: $\exists a, b \in \R : \bigl(a < b \land \forall q \in A : (q \le a \lor q \ge b)\bigr)$. Some open interval misses $A$. For $\Z$: $a = 0$, $b = 1$.

$\Q$ is dense: given $a < b$, choose $n \in \N$ with $1/n < b - a$ (Archimedean property). Let $m$ be the least integer with $m > na$. Then $m/n > a$, and by minimality $m - 1 \le na$, so $m/n \le a + 1/n < b$.

\ans{ex:open}
Not open:
\[
  \exists x \in U\ \forall r > 0\ \exists y \in \R : \bigl(|y - x| < r \land y \notin U\bigr).
\]
Some point of $U$ has points outside $U$ arbitrarily close to it.

$[0,1)$ is not open: take $x = 0$; given $r$, take $y = -r/2$. It is not closed: the complement $C = (-\infty, 0) \cup [1, \infty)$ is not open, with witness $x = 1 \in C$; given $r$, take $y = 1 - \min(r, 1)/2$, which lies in $[0,1)$ and satisfies $|y - 1| < r$.

$\emptyset$ is open because the statement $\forall x \in \emptyset \dots$ is vacuously true. Its complement $\R$ is open with $r = 1$ for every $x$. So $\emptyset$ is both open and closed, which confirms that ``closed'' is not the negation of ``open.''

\ans{ex:limpt}
Not a limit point: $\exists \delta > 0\ \forall a \in A : (a = c \lor |a - c| \ge \delta)$. Some neighborhood of $c$ contains no point of $A$ other than possibly $c$ itself.

For $c \in \Z$: $\delta = 1/2$ works, since distinct integers are at distance at least $1$. For $\{1/n\}$ and $c = 0$: given $\delta$, choose $n > 1/\delta$; then $0 < 1/n < \delta$. And $0 \ne 1/n$ for all $n$.

\ans{ex:seqcpt}
Sequentially compact:
\[
  \forall (x_n) \subseteq K\ \exists (n_k) \text{ strictly increasing}\ \exists x \in K : x_{n_k} \to x.
\]
Negation:
\[
  \exists (x_n) \subseteq K\ \forall (n_k) \text{ strictly increasing}\ \forall x \in K : x_{n_k} \not\to x.
\]
There is a sequence in $K$ no subsequence of which converges to a point of $K$.

$(0,1]$: take $x_n = 1/n$. Every subsequence converges to $0$, limits are unique, and $0 \notin K$. So no subsequence converges to a point of $K$. The failure is a missing limit point: $K$ is not closed.

$\R$: take $x_n = n$. Every subsequence is unbounded, and convergent sequences are bounded. The failure is unboundedness. By the Bolzano--Weierstrass theorem, these are the only two ways a subset of $\R$ can fail: it is sequentially compact if and only if it is closed and bounded.

\subsection*{Synthesis}

\ans{ex:equi}
Negation:
\[
  \exists \epsilon > 0\ \forall \delta > 0\ \exists f \in \mathcal{F}\ \exists x, y \in D : \bigl(|x - y| < \delta \land |f(x) - f(y)| \ge \epsilon\bigr).
\]
For $\{x^n\}$ on $[0,1]$: take $\epsilon = 1/2$. Given $\delta$, let $y = 1$ and $x = \max(0, 1 - \delta/2)$. Then $0 \le x < 1$ and $|x - y| \le \delta/2 < \delta$. Since $x^n \to 0$, choose $n$ with $x^n < 1/2$. Then $|f(x) - f(y)| = 1 - x^n > 1/2$ for $f(t) = t^n$. Each $t^n$ is uniformly continuous on $[0,1]$ (it is continuous on a closed bounded interval), but its $\delta$ must shrink as $n$ grows, and no single $\delta$ serves the whole family. Compare Exercise~\ref{ex:ucimpl}: equicontinuity moves $\exists \delta$ in front of $\forall f$, just as uniform continuity moves it in front of $\forall c$.

\ans{ex:everycont}
Statement:
\[
  \forall f \colon D \to \R : \Bigl( f \text{ continuous on } D \implies \exists M\ \forall x \in D : |f(x)| \le M \Bigr).
\]
Negation:
\[
  \exists f \colon D \to \R : \Bigl( f \text{ continuous on } D \land \forall M\ \exists x \in D : |f(x)| > M \Bigr).
\]
For $D = (0,1)$, take $f(x) = 1/x$. Given $M$, let $M' = \max(M, 1)$ and $x = 1/(2M')$. Then $x \in (0, 1/2]$ and $f(x) = 2M' > M$. For $D = [0,1]$, the extreme value theorem says every continuous $f$ is bounded and attains its bounds, so the statement is true and the negation false. The negation of a theorem is a precise description of what a counterexample must be.

\ans{ex:swapgeneral}
Suppose $y_0$ satisfies $\forall x\ P(x, y_0)$. Then for each $x$, $y = y_0$ is a witness for $\exists y\ P(x, y)$. Counterexample to the converse: over $\R$, $P(x, y) \equiv y > x$ (Exercise~\ref{ex:swap1}).

Same-type quantifiers commute: $\forall x\ \forall y\ P$ and $\forall y\ \forall x\ P$ both say $P$ holds for every pair, and $\exists x\ \exists y\ P$ and $\exists y\ \exists x\ P$ both say some pair satisfies $P$.

Instances of $\exists\forall \implies \forall\exists$, each strict: uniform continuity implies continuity (Exercise~\ref{ex:ucimpl}); uniform convergence implies pointwise convergence (Exercise~\ref{ex:pw}); equicontinuity of a family implies that each member is uniformly continuous (Exercise~\ref{ex:equi}). Exercise~\ref{ex:swapconv} is the same move applied to convergence itself, where it collapses the definition to eventual constancy.

\end{document}
